\documentclass{beamer}
\usetheme{Madrid}
\usecolortheme{default}
\usefonttheme{professionalfonts}

\usepackage{graphicx}
\usepackage{amsmath}
\usepackage{bm}        % bold vectors/tensors
\usepackage{hyperref}  % robust \href

\title{The Hybrid Immersed Boundary Finite Element Method Introduction}
\author{Laura Miller}
\date{\today}

\begin{document}

\frame{\titlepage}

\begin{frame}{Resources}
  \begin{itemize}
    \item \href{https://www.youtube.com/watch?v=-1TlNJL93sc}{Boyce Griffith presents IBFE at 23-minute mark.}
    \item \href{https://drive.google.com/drive/folders/1vkhrX0YInMm5pfUj3LcRA9TV0DZdFp1a?usp=sharing}{Matea presents solid mechanics and IBFE.}
    \item \href{https://www.classcentral.com/classroom/youtube-aim-micde-seminar-boyce-griffith-immersed-methods-for-fluid-structure-interaction-323154}{Immersed Methods for FSI by Boyce Griffith}
  \end{itemize}
\end{frame}

\begin{frame}{Why IBFE? The core difficulty in fluid--structure interaction}
\small
In FSI we need to evolve:
\begin{itemize}
  \item an \textbf{incompressible viscous fluid} on an Eulerian grid (fixed mesh),
  \item a \textbf{deformable solid} that can move and undergo large strains,
  \item and enforce \textbf{kinematic coupling} (no-slip) and \textbf{dynamic coupling} (force balance) at the interface.
\end{itemize}

\medskip
\textbf{Main numerical bottleneck:} the structure can move a lot, change shape, and even self-contact.
Body-fitted meshes must be updated or remeshed, which can be expensive and fragile.

\medskip
\textbf{IBFE idea:} keep the \emph{fluid grid fixed} (Eulerian) and represent the \emph{solid mechanics} with a \emph{Lagrangian finite element mesh}.
Couple them with smooth integral operators.
\end{frame}


\begin{frame}{FSI strategies: where IBFE fits}
\small
\begin{itemize}
  \item \textbf{Body-fitted / ALE methods:}
  mesh conforms to the interface. Very accurate at boundaries, but large deformation can require remeshing.
  \item \textbf{Classic Immersed Boundary (Peskin IBM):}
  structure represented by Lagrangian markers/springs; forces spread to Eulerian grid with a regularized delta.
  Great for moving boundaries, but \emph{structural model can be low-order} unless enhanced.
  \item \textbf{IBFE (hybrid IB + FEM):}
  keep the immersed coupling of IBM \emph{and} use FE to compute solid stress/forces with standard hyperelasticity.
\end{itemize}

\medskip
\textbf{Takeaway:} IBFE is ``immersed coupling'' + ``finite-element solid mechanics''.
\end{frame}

\begin{frame}{What IBFE adds beyond classic IBM}
\small
Classic IBM often models the structure with springs/fibers/penalty forces on markers.
IBFE replaces that piece with \textbf{continuum solid mechanics}:

\medskip
\begin{itemize}
  \item \textbf{Constitutive modeling:} hyperelasticity/viscoelasticity via \(W(\bm{F})\), stresses \( \bm{P}=\partial W/\partial \bm{F}\).
  \item \textbf{Accuracy/robustness:} FE representation supports higher-order elements, anisotropy, nearly incompressible solids, etc.
  \item \textbf{Clean variational formulation:} internal forces come from weak form (stable and standard in FEM).
\end{itemize}

\medskip
\textbf{What does not change:} the coupling operators (spread/interpolate) and the fixed Eulerian fluid grid.
\end{frame}

\begin{frame}{IBFE in one sentence (and one picture)}
\small
\textbf{IBFE = Eulerian fluid solver + Lagrangian FE solid + smooth coupling operators.}

\medskip
\begin{itemize}
  \item Fluid: incompressible Navier--Stokes on a fixed grid (efficient and mature).
  \item Solid: finite-strain continuum mechanics on a moving FE mesh (accurate material modeling).
  \item Coupling: regularized delta / kernel operators that transfer \(\bm{u}\) and \(\bm{f}\) between meshes.
\end{itemize}

\medskip
\textbf{Interpretation:} the FE mesh is \emph{not} used to remesh the fluid.
It is an immersed ``computational skeleton'' that carries solid kinematics and constitutive laws.
\end{frame}

\begin{frame}{IBFE pipeline: what gets computed each time step?}
\small
At time \(t^n\), we have the structure configuration \(\chi(\cdot,t^n)\) and fluid velocity \(\bm{u}(\cdot,t^n)\).

\medskip
\begin{enumerate}
  \item \textbf{Interpolate} Eulerian velocity to the Lagrangian mesh:
  \(\ \partial_t\chi(\bm{X},t^n) = \int_{\Omega} \bm{u}(\bm{x},t^n)\,\delta_h(\bm{x}-\chi(\bm{X},t^n))\,d\bm{x}\).
  \item \textbf{Compute} solid deformation measures on the FE mesh:
  \(\bm{F}=\nabla_X\chi,\ J=\det\bm{F}\), etc.
  \item \textbf{Compute} solid stress from a constitutive law:
  \(\bm{P}=\partial W/\partial \bm{F}\) (plus viscosity/active terms if used).
  \item \textbf{Assemble} the Lagrangian force density (weak form \(\Rightarrow\) nodal forces \(\Rightarrow\) \(\bm{F}(\bm{X},t^n)\)).
  \item \textbf{Spread} Lagrangian force density to the Eulerian grid:
  \(\ \bm{f}(\bm{x},t^n)=\int \bm{F}(\bm{X},t^n)\,\delta_h(\bm{x}-\chi(\bm{X},t^n))\,d\bm{X}\).
  \item \textbf{Advance} the fluid (Navier--Stokes with body force \(\bm{f}\)).
\end{enumerate}

\medskip
\textbf{Takeaway:} FE is used to compute \(\bm{P}\) and internal forces; the fluid is still solved on a fixed Eulerian grid.
\end{frame}


\begin{frame}{Coupling the Methods}
    \begin{itemize}
        \item The coupling between the immersed boundary method and the finite element method is achieved through integral equations that relate the Lagrangian and Eulerian variables.
        \item These integral equations involve:
        \begin{itemize}
            \item Force spreading: Distributing the Lagrangian forces from the immersed structure to the Eulerian fluid grid.
            \item Velocity interpolation: Interpolating the Eulerian fluid velocities to the Lagrangian structure.
        \end{itemize}
        \item The integral equations ensure a smooth and accurate transfer of information between the two discretizations.
    \end{itemize}
\end{frame}


\begin{frame}{Forces in immersed methods: what is actually coupled?}
\small
Immersed methods couple the fluid and structure through a \textbf{force density} added to the fluid momentum equation:
\[
\rho\left(\partial_t\bm{u}+\bm{u}\cdot\nabla\bm{u}\right) = -\nabla p + \mu\Delta \bm{u} + \bm{f},
\qquad \nabla\cdot\bm{u}=0.
\]

\medskip
\textbf{In IB/IBFE:}
\begin{itemize}
  \item The fluid sees an \textbf{Eulerian body force density} \(\bm{f}(\bm{x},t)\) (units: force per volume).
  \item That force comes from a \textbf{Lagrangian force density} \(\bm{F}(\bm{X},t)\) computed from solid stress/elasticity (plus any active terms).
\end{itemize}

\medskip
\textbf{Traction viewpoint:} when needed, surface tractions satisfy \(\bm{t}=\boldsymbol{\sigma}\bm{n}\),
but the numerical coupling is usually implemented via volumetric force densities and weak forms.
\end{frame}


\begin{frame}{Traction Forces}
    \begin{itemize}
        \item Traction forces are generated by the fluid's stress acting on the surface of the immersed structure.
        \item The traction vector $\mathbf{t}$ at a point on the surface is given by:
        \begin{equation*}
        \mathbf{t} = \boldsymbol{\sigma} \mathbf{n}
        \end{equation*}
        where:
        \begin{itemize}
            \item $\boldsymbol{\sigma}$ is the Cauchy stress tensor in the fluid
            \item $\mathbf{n}$ is the outward unit normal vector to the surface
        \end{itemize}
        \item The traction force can be decomposed into normal and tangential components:
        \begin{itemize}
            \item Normal traction: Acts perpendicular to the surface.
            \item Tangential traction (shear stress): Acts parallel to the surface.
        \end{itemize}
    \end{itemize}
\end{frame}

\begin{frame}{Body Forces}
    \begin{itemize}
        \item Body forces act on the volume of the immersed structure.
        \item Examples of body forces include:
        \begin{itemize}
            \item Gravity: Acts on the mass of the structure.
            \item Active tension: Generated by muscle contraction or other active elements within the structure.
            \item External forces: Applied to the structure from outside the fluid domain.
        \end{itemize}
        \item Body forces are incorporated into the equations of motion for the structure.
    \end{itemize}
\end{frame}

\begin{frame}{Eulerian and Lagrangian Equations of Motion}
\begin{itemize}
\item In this formulation, the structural kinematics are described in Lagrangian form using the mapping $\chi: (\Omega_s^0, t) \to \Omega_s^t$, which connects the reference configuration to the current configuration at time $t$. 
\item Specifically, $\chi(\mathbf{X}, t) \in \Omega_s^t$ is the position at time $t$ of material point $\mathbf{X} \in \Omega_s^0$. 
\item A composite material velocity field is introduced on $\Omega$ so that $\mathbf{u}(\mathbf{x}, t)$ is the velocity of whichever material happens to be located at position $\mathbf{x}$ at time $t$. 
\end{itemize}
\end{frame}

\begin{frame}{Stress and Strain}
    \begin{figure}
        \includegraphics[width=0.9\textwidth]{pictures/lagrangian-eulerian-coordinate.png}
        \caption{Lagrangian and Eulerian coordinate systems. The Lagrangian material coordinate domain is $U$, and the Eulerian physical coordinate domain is $\Omega$. The physical position of material point $X$ at time $t$ is $\chi(X, t)$, the physical region occupied by the structure is $\chi(U,t)$, and the physical region occupied by the fluid is $\Omega \setminus \chi(U,t)$.}
    \end{figure}
\end{frame}

\begin{frame}{Why does IBFE use the first Piola--Kirchhoff stress $\bm{P}$?}
\small
The FE mesh lives in \textbf{reference coordinates} \(\bm{X}\in \Omega_s^0\).
Internal forces are naturally written per \textbf{reference area/volume}.

\medskip
\textbf{Hyperelasticity:} choose a strain energy density \(W(\bm{F})\) per reference volume, then
\[
\bm{P}(\bm{X},t) = \frac{\partial W}{\partial \bm{F}}(\bm{X},t).
\]

\medskip
\textbf{Bridge to spatial stress (if you want it):}
\[
\bm{P} = J\,\boldsymbol{\sigma}\,\bm{F}^{-T},
\qquad J=\det\bm{F}.
\]

\medskip
\textbf{IBFE takeaway:} we compute \(\bm{P}\) on the FE mesh, then convert it into nodal forces via the weak form.
\end{frame}


\begin{frame}{First Piola-Kirchhoff Stress Tensor}
\begin{itemize}
\item The formulation uses the first Piola-Kirchhoff solid stress tensor $\mathbf{P}_s$, which is related to the corresponding Cauchy stress tensor $\boldsymbol{\sigma}_s$ by: 
\begin{equation*}
\mathbf{P}_s(\mathbf{X}, t) = J(\mathbf{X}, t) \boldsymbol{\sigma}_s[\chi(\mathbf{X}, t), t] \mathbf{F}(\mathbf{X}, t)^{-T}
\end{equation*}
where:
\begin{itemize}
\item $\mathbf{F} = \frac{\partial \chi}{\partial \mathbf{X}}$ is the deformation gradient tensor. 
\item $J = \det(\mathbf{F})$ is the Jacobian determinant. 
\end{itemize}
\end{itemize}
\end{frame}

\begin{frame}{Assumptions}
\begin{itemize}
\item The discussion is restricted to a neutrally buoyant, incompressible structure immersed in a viscous incompressible fluid, so that $\rho$ is the mass density of both the fluid and the solid. 
\item The structure is assumed to be viscoelastic, and the form of the solid viscosity is the same as that of the fluid. 
\end{itemize}
\end{frame}



\begin{frame}{Coupling Equations in Detail}
    \begin{itemize}
        \item The integral form of the coupling equations ensures a smooth and accurate transfer of information between the Lagrangian and Eulerian discretizations.
        \item The Dirac delta function acts as a kernel that spreads the Lagrangian forces to the Eulerian grid and interpolates the Eulerian velocities to the Lagrangian structure.
        \item The choice of the Dirac delta function or other interpolation kernels can affect the accuracy and stability of the method.
    \end{itemize}
\end{frame}

\begin{frame}{Elasticity Models in the Immersed Boundary Method}
\begin{itemize}
\item The immersed boundary method can incorporate various elasticity models to describe the behavior of the immersed structure.
\item These models typically involve a strain energy density function that relates the deformation of the structure to the stored elastic energy.
\item The choice of elasticity model depends on the specific material properties and the type of deformation the structure undergoes.
\end{itemize}
\end{frame}

\begin{frame}{Hyperelasticity}
\begin{itemize}
\item Hyperelastic materials are a class of materials that exhibit nonlinear elastic behavior and can undergo large deformations.
\item Their constitutive behavior is described by a strain energy density function $W$ that depends on the deformation gradient tensor $\mathbf{F}$.
\item The first Piola-Kirchhoff stress tensor $\mathbf{P}$ is given by the derivative of the strain energy density function with respect to the deformation gradient tensor:
\begin{equation*}
\mathbf{P} = \frac{\partial W}{\partial \mathbf{F}}
\end{equation*}
\end{itemize}
\end{frame}



\begin{frame}{Applying Active Tension}
    \begin{itemize}
        \item Active tension refers to the force generated by the active contraction of muscles or other actuators within the structure.
        \item In the hybrid method, active tension can be incorporated by adding an active force term to the Lagrangian force density $\mathbf{F}$.
        \item The active tension can be modeled using various constitutive laws or prescribed directly based on experimental data.
    \end{itemize}
\end{frame}

\begin{frame}{Advantages of the Hybrid Approach}
    \begin{itemize}
        \item The hybrid method offers several advantages, including:
        \begin{itemize}
            \item Ability to handle complex geometries and large deformations
            \item Improved accuracy and efficiency compared to the original immersed boundary method
            \item Versatility in handling a wide range of material properties and boundary conditions
        \end{itemize}
    \end{itemize}
\end{frame}

\begin{frame}{Applications of the Hybrid Method}
    \begin{itemize}
        \item The hybrid method has been successfully applied to various fluid-structure interaction problems, such as:
        \begin{itemize}
            \item Simulation of blood flow in the heart
            \item Flapping of insect wings
            \item Deformation of soft tissues
            \item Interaction of fluids with flexible structures
        \end{itemize}
    \end{itemize}
\end{frame}

\begin{frame}{Limitations and Future Directions}
    \begin{itemize}
        \item Despite its advantages, the hybrid method has some limitations, such as:
        \begin{itemize}
            \item Computational cost can be high for complex problems
            \item Numerical stability can be challenging to maintain for certain parameter regimes
        \end{itemize}
        \item Future research directions include:
        \begin{itemize}
            \item Developing more efficient and stable numerical algorithms
            \item Extending the method to handle more complex material properties and boundary conditions
            \item Applying the method to new and challenging fluid-structure interaction problems
        \end{itemize}
    \end{itemize}
\end{frame}

\begin{frame}{When to use IBFE (and when not to)}
\small
\textbf{IBFE shines when:}
\begin{itemize}
  \item large deformation or moving boundaries make body-fitted meshing painful,
  \item you need serious solid mechanics (anisotropy, nearly incompressible solids, active stress),
  \item topology stays fixed (no tearing/merging of the solid mesh).
\end{itemize}

\medskip
\textbf{Tradeoffs:}
\begin{itemize}
  \item Interface is \emph{regularized} by the kernel \(\delta_h\): sharp boundary layers are smeared over a few grid cells.
  \item Nearly incompressible solids can be stiff \(\Rightarrow\) time-step restrictions or need for implicit/semi-implicit treatment.
  \item Contact/self-contact needs special handling (as in many FSI methods).
\end{itemize}

\medskip
\textbf{Rule of thumb:}
IBFE is a great default for soft-tissue-type FSI where geometry moves a lot and constitutive modeling matters.
\end{frame}


\begin{frame}{Conclusion}
    \begin{itemize}
        \item The hybrid immersed boundary finite element method is a powerful tool for simulating complex fluid-structure interaction problems.
        \item It combines the strengths of both the immersed boundary method and the finite element method, offering improved accuracy, efficiency, and versatility.
        \item This method has broad applications in various fields and continues to be an active area of research and development.
    \end{itemize}
\end{frame}

\begin{frame}{Glossary of Terms}
  \begin{description}
    \item[Stress] A measure of the force acting per unit area within a deformable body.
    \item[Tensor] A mathematical object that describes a multilinear relationship between sets of algebraic objects related to a vector space.
    \item[Cauchy Stress Tensor ($\boldsymbol{\sigma}$)] A second-order tensor that describes the state of stress at a point within a material.
    \item[Normal Stress] Stress acting perpendicular to a surface.
    \item[Shear Stress] Stress acting parallel to a surface.
    \item[Pascal (Pa)] The SI unit of pressure or stress, equal to one newton per square meter.
    \item[Pound per Square Inch (psi)] A unit of pressure or stress in the imperial and US customary systems.
    \item[Continuum Mechanics] The branch of mechanics that deals with the analysis of the kinematics and mechanical behavior of materials modeled as a continuous mass rather than as discrete particles.
  \end{description}
\end{frame}

\end{document}